\documentclass[11pt]{article}

\usepackage[margin=1in]{geometry}
\usepackage[T1]{fontenc}
\usepackage[utf8]{inputenc}
\usepackage{lmodern}
\usepackage{amsmath,amssymb,amsfonts,bm}
\usepackage{mathtools}
\usepackage{graphicx}
\usepackage{booktabs}
\usepackage{xcolor}
\usepackage{hyperref}
\usepackage{enumitem}
\usepackage{siunitx}

\hypersetup{
    colorlinks=true,
    linkcolor=blue,
    citecolor=blue,
    urlcolor=blue
}

\title{\textbf{Study Plan: Predictive Modeling of Measurement-Induced Leakage and Transmon Ionization at High Readout Power}\\
\vspace{0.3em}
\large A \texttt{cqed\_sim}-Based Numerical and Experimental Validation Program}
\author{}
\date{\today}

\begin{document}

\maketitle

\section{Project Title}

\textbf{Predictive Modeling of Measurement-Induced Leakage and Transmon Ionization at High Readout Power}

\section{Problem Statement}

The goal of this project is to develop a quantitative simulation framework in
\texttt{cqed\_sim} that predicts the onset and severity of measurement-induced transitions
out of the computational manifold during high-power dispersive readout of a transmon-cavity
system. In particular, the study will model population transfer into non-computational
transmon levels, including $|f\rangle$ and higher-lying states, and characterize the
regime in which strong measurement drives produce effective ionization-like behavior into
high excited manifolds.

The simulator should predict leakage and non-QND behavior as a function of:
\begin{itemize}[leftmargin=2em]
    \item readout drive amplitude,
    \item readout drive detuning,
    \item resonator linewidth $\kappa$,
    \item dispersive shift $\chi$,
    \item transmon anharmonicity,
    \item qubit-cavity detuning and coupling,
    \item cavity and transmon nonlinearities,
    \item pulse shape and readout-window duration,
    \item effective filter-chain bandwidth and transfer-function distortions.
\end{itemize}

The final goal is not merely to simulate assignment fidelity, but to build a predictive
tool that can identify the onset of QND breakdown, quantify leakage into higher levels, and
guide the experimental design of faster and safer readout protocols.

\section{Why This Matters Experimentally}

Fast, high-fidelity readout is central to repeated syndrome extraction, real-time feedback,
active reset, and quantum error correction. However, pushing readout toward shorter windows
and larger cavity photon number can invalidate the weak-drive dispersive approximation. In
this regime, the measurement tone itself can induce transitions between dressed states,
cause leakage into $|f\rangle$ and higher transmon levels, and eventually populate highly
excited states in a manner that is operationally ionization-like.

These effects matter because they:
\begin{itemize}[leftmargin=2em]
    \item break QND assumptions,
    \item introduce state mixing during measurement,
    \item reduce repeated-measurement consistency,
    \item harm active reset and feedback fidelity,
    \item complicate interpretation of integrated homodyne signals,
    \item create hidden error channels relevant for QEC cycles.
\end{itemize}

A predictive simulator would therefore be directly useful both scientifically and
operationally: it would allow one to identify safe operating regions, understand breakdown
mechanisms, and optimize the tradeoff between speed, fidelity, and leakage.

\section{Scientific Objectives}

This study should answer the following questions:

\begin{enumerate}[leftmargin=2em]
    \item At what drive amplitude, pulse duration, and detuning does readout cease to be
    approximately QND?
    \item How do leakage thresholds scale with $\kappa$, $\chi$, anharmonicity, and pulse
    shaping?
    \item When is a simple weakly anharmonic Duffing approximation insufficient, and when
    must one use a more accurate multi-level transmon basis?
    \item Can one distinguish moderate leakage into $|f\rangle$ from more severe
    high-lying-state population spreading?
    \item How do finite-bandwidth filtering and pulse distortion alter the leakage
    threshold?
    \item Which mitigation strategies best preserve fidelity while suppressing leakage and
    residual photons?
\end{enumerate}

\section{Physical Model}

\subsection{Baseline Hamiltonian}

The core system is a driven cavity coupled to a multi-level transmon. The simulation should
be performed in an appropriate rotating frame. The time-dependent Hamiltonian may be written
schematically as
\begin{equation}
    H(t) = H_{\mathrm{sys}} + H_{\mathrm{drive}}(t),
\end{equation}
where $H_{\mathrm{sys}}$ includes the bare cavity, the multi-level transmon, their
interaction, and any relevant nonlinear terms, while $H_{\mathrm{drive}}(t)$ describes the
measurement drive applied through the readout channel.

A minimal model should include:
\begin{itemize}[leftmargin=2em]
    \item a single readout cavity mode,
    \item a multi-level transmon,
    \item qubit-cavity coupling,
    \item coherent measurement drive,
    \item cavity decay,
    \item transmon relaxation and dephasing.
\end{itemize}

\subsection{Model Extensions for Strong Readout}

Because the project targets the high-power regime, the study should allow the following
extensions:
\begin{itemize}[leftmargin=2em]
    \item cavity self-Kerr,
    \item inherited qubit-induced cavity nonlinearity,
    \item realistic multi-level transmon matrix elements,
    \item filtered and distorted time-domain drive envelopes,
    \item optional effective Floquet treatment for strong periodic drive,
    \item optional polaron-based reduced descriptions in the dispersive regime.
\end{itemize}

\subsection{Recommended Transmon Modeling Hierarchy}

The numerical study should proceed in a layered way:
\begin{enumerate}[leftmargin=2em]
    \item \textbf{Fast coarse scans:} weakly anharmonic oscillator or Duffing transmon
    model.
    \item \textbf{Quantitative leakage prediction:} multi-level transmon basis built from
    accurate eigenstates and operator matrix elements.
    \item \textbf{High-power regime refinement:} larger transmon truncation and careful
    convergence tests near the apparent ionization boundary.
\end{enumerate}

\section{Open-System Dynamics}

\subsection{Unconditional Evolution}

For unconditional dynamics, the simulator should solve a Lindblad master equation of the
form
\begin{equation}
    \dot{\rho} = -i[H(t), \rho] + \sum_j \mathcal{D}[L_j]\rho,
\end{equation}
with
\begin{equation}
    \mathcal{D}[L]\rho = L\rho L^\dagger - \frac{1}{2}\left(L^\dagger L \rho + \rho L^\dagger L\right).
\end{equation}

Relevant collapse operators should include:
\begin{itemize}[leftmargin=2em]
    \item cavity decay,
    \item transmon energy relaxation,
    \item transmon dephasing,
    \item optional thermal excitation,
    \item optional effective incoherent processes if justified by the experiment.
\end{itemize}

\subsection{Conditioned Evolution}

To connect the simulation to real readout records, the framework should also support
stochastic master equations or quantum trajectories. This is necessary to generate:
\begin{itemize}[leftmargin=2em]
    \item synthetic homodyne or heterodyne measurement records,
    \item integrated readout histograms,
    \item assignment fidelities,
    \item shot-to-shot leakage statistics,
    \item jump-like trajectory signatures associated with non-QND transitions.
\end{itemize}

\section{Core Observables}

The simulator should extract the following observables:
\begin{itemize}[leftmargin=2em]
    \item transmon level populations $P_n(t)$,
    \item computational-subspace populations $P_g(t)$ and $P_e(t)$,
    \item $|f\rangle$ population $P_f(t)$,
    \item leakage fraction
    \begin{equation}
        P_{\mathrm{leak}}(t)=1-P_g(t)-P_e(t),
    \end{equation}
    \item cavity photon number $\langle a^\dagger a \rangle(t)$,
    \item peak intracavity photon number during measurement,
    \item post-pulse residual cavity population,
    \item assignment fidelity from simulated readout records,
    \item repeated-measurement transition probabilities,
    \item broadening, skewness, and non-Gaussian features of the integrated homodyne
    distributions.
\end{itemize}

\subsection{Operational Ionization Metric}

Because true ionization is not literal in a finite truncated Hilbert space, the report
should define an operational high-lying-state metric. For a chosen cutoff
$n_{\mathrm{ion}}$, define
\begin{equation}
    P_{\mathrm{ion}}(t)=\sum_{m\ge n_{\mathrm{ion}}} P_m(t).
\end{equation}

This metric should be used to distinguish:
\begin{itemize}[leftmargin=2em]
    \item safe dispersive readout,
    \item moderate leakage into nearby non-computational levels,
    \item severe population spreading into high-lying transmon states.
\end{itemize}

\section{Parameter Regimes and Sweep Axes}

\subsection{Drive Parameters}
\begin{itemize}[leftmargin=2em]
    \item readout amplitude,
    \item readout detuning from the dressed cavity resonance,
    \item pulse length,
    \item rise and fall times,
    \item pulse-envelope shape,
    \item optional two-tone or shelving-assisted variants.
\end{itemize}

\subsection{Device Parameters}
\begin{itemize}[leftmargin=2em]
    \item $\kappa$,
    \item $\chi$,
    \item transmon anharmonicity,
    \item qubit-cavity detuning,
    \item effective coupling strength,
    \item cavity Kerr and induced nonlinearities.
\end{itemize}

\subsection{Control-Chain / Filter Parameters}
\begin{itemize}[leftmargin=2em]
    \item finite analog bandwidth,
    \item low-pass filtering,
    \item rise-time limitation,
    \item overshoot or ringing,
    \item effective transfer-function distortion of the intended pulse.
\end{itemize}

\subsection{Numerical Convergence Parameters}
\begin{itemize}[leftmargin=2em]
    \item number of transmon levels,
    \item cavity photon truncation,
    \item solver time-step and tolerances,
    \item number of quantum trajectories.
\end{itemize}

\section{Work Packages}

\subsection{WP1: Baseline Simulator Assembly}

\textbf{Objective.} Build a robust driven multi-level transmon plus cavity simulation module
inside \texttt{cqed\_sim}.

\textbf{Tasks.}
\begin{itemize}[leftmargin=2em]
    \item Implement Hamiltonian construction for cavity plus multi-level transmon.
    \item Implement collapse-operator generation.
    \item Implement pulse-envelope and transfer-function modules.
    \item Implement observable extraction and data logging.
    \item Build a parameterized sweep framework.
\end{itemize}

\textbf{Deliverables.}
\begin{itemize}[leftmargin=2em]
    \item solver-ready model builder,
    \item reproducible config files,
    \item automated sweep scripts,
    \item baseline figures verifying expected weak-drive behavior.
\end{itemize}

\subsection{WP2: Leakage Threshold Mapping}

\textbf{Objective.} Identify the onset of measurement-induced leakage and non-QND behavior
across operating regimes.

\textbf{Tasks.}
\begin{itemize}[leftmargin=2em]
    \item Sweep readout amplitude and pulse duration.
    \item Sweep detuning around the dressed cavity resonance.
    \item Sweep $\kappa$, $\chi$, and anharmonicity.
    \item Construct leakage and ionization heatmaps.
\end{itemize}

\textbf{Key Outputs.}
\begin{itemize}[leftmargin=2em]
    \item $P_f$ after readout,
    \item total leakage $P_{\mathrm{leak}}$,
    \item high-level population $P_{\mathrm{ion}}$,
    \item fidelity versus leakage tradeoff curves,
    \item residual-photon versus speed tradeoff curves.
\end{itemize}

\subsection{WP3: Conditioned Readout Modeling}

\textbf{Objective.} Connect physical leakage mechanisms to experimentally visible readout
signatures.

\textbf{Tasks.}
\begin{itemize}[leftmargin=2em]
    \item Generate synthetic homodyne trajectories.
    \item Integrate the readout record using realistic kernels.
    \item Construct assignment histograms and confusion matrices.
    \item Identify signatures of jump-like state changes under strong drive.
\end{itemize}

\textbf{Key Outputs.}
\begin{itemize}[leftmargin=2em]
    \item readout histograms below, near, and above threshold,
    \item broadened or non-Gaussian integrated distributions,
    \item state-assignment errors attributable to true state mixing rather than pure
    amplifier noise.
\end{itemize}

\subsection{WP4: High-Level-State / Ionization Study}

\textbf{Objective.} Characterize the regime where strong drive populates high transmon
manifolds.

\textbf{Tasks.}
\begin{itemize}[leftmargin=2em]
    \item increase transmon truncation,
    \item monitor population spreading across excited manifolds,
    \item compare reduced models to more accurate multi-level representations,
    \item identify practical boundaries between ordinary leakage and severe high-level
    occupation.
\end{itemize}

\textbf{Success Criterion.}
A robust regime map separating:
\begin{itemize}[leftmargin=2em]
    \item approximately QND readout,
    \item moderate leakage,
    \item strong non-QND mixing,
    \item high-lying-state or ionization-like behavior.
\end{itemize}

\subsection{WP5: Mitigation and Control Study}

\textbf{Objective.} Evaluate practical mitigation strategies at fixed fidelity targets.

\textbf{Candidate Strategies.}
\begin{itemize}[leftmargin=2em]
    \item pulse shaping to reduce peak photon number,
    \item smoother ring-up and ring-down envelopes,
    \item detuning optimization,
    \item two-tone or shelving-assisted readout,
    \item $\chi/\kappa$ redesign strategies,
    \item filter-chain optimization,
    \item reduced residual-photon protocols.
\end{itemize}

\textbf{Output.}
A Pareto analysis over:
\begin{itemize}[leftmargin=2em]
    \item assignment fidelity,
    \item leakage,
    \item high-lying-state population,
    \item pulse duration,
    \item post-readout residual cavity population.
\end{itemize}

\section{Validation Experiments}

\subsection{Purpose of Validation}

The simulator must be validated against experimentally measurable observables that directly
probe non-QND effects, not just assignment fidelity. The validation goal is to show that
the model predicts both:
\begin{enumerate}[leftmargin=2em]
    \item the onset of leakage and state mixing,
    \item the structure of observable measurement signatures associated with that leakage.
\end{enumerate}

\subsection{Experiment A: Leakage vs Readout Power}

\textbf{Protocol.}
\begin{itemize}[leftmargin=2em]
    \item Prepare $|g\rangle$ and $|e\rangle$.
    \item Apply a readout pulse with fixed shape and duration.
    \item Sweep readout power from the low-power dispersive regime into the strong-drive
    regime.
    \item Perform post-readout level discrimination, ideally resolving at least
    $|g\rangle$, $|e\rangle$, and $|f\rangle$.
\end{itemize}

\textbf{Measured observables.}
\begin{itemize}[leftmargin=2em]
    \item post-readout $P_g$, $P_e$, $P_f$,
    \item total leakage fraction,
    \item assignment fidelity.
\end{itemize}

\textbf{Validation target.}
Compare the simulated and measured onset power for leakage, the slope of leakage growth,
and the dependence on the initial state.

\subsection{Experiment B: Leakage vs Pulse Duration}

\textbf{Protocol.}
\begin{itemize}[leftmargin=2em]
    \item Fix the nominal readout amplitude.
    \item Sweep pulse duration and rise/fall time.
    \item Measure level populations after the readout pulse.
\end{itemize}

\textbf{Validation target.}
Determine whether the simulation correctly captures the crossover between adiabatic and
non-adiabatic regimes, and whether leakage is driven mainly by peak power or by total
exposure time.

\subsection{Experiment C: Detuning Dependence}

\textbf{Protocol.}
\begin{itemize}[leftmargin=2em]
    \item Sweep readout tone detuning around the dressed cavity resonance.
    \item Measure post-readout leakage and assignment fidelity.
\end{itemize}

\textbf{Validation target.}
Test whether the simulator correctly identifies dangerous detuning regions where level
hybridization or unwanted transitions are enhanced.

\subsection{Experiment D: Repeated-Measurement QND Test}

\textbf{Protocol.}
\begin{itemize}[leftmargin=2em]
    \item Prepare a known state.
    \item Perform one readout.
    \item Optionally wait a controlled delay.
    \item Perform a second readout.
    \item Compare the second outcome conditioned on the first.
\end{itemize}

\textbf{Measured observables.}
\begin{itemize}[leftmargin=2em]
    \item repeated-measurement consistency,
    \item conditional transition probabilities,
    \item apparent readout-induced excitation or relaxation rates.
\end{itemize}

\textbf{Validation target.}
Compare simulated and experimental QND breakdown thresholds and transition rates.

\subsection{Experiment E: Higher-Level Resolution via Shelving or Two-Tone Probing}

\textbf{Protocol.}
\begin{itemize}[leftmargin=2em]
    \item Use a post-readout spectroscopy or shelving sequence to discriminate $|f\rangle$
    and, where feasible, selected higher excited states.
    \item Repeat across readout power and pulse-length sweeps.
\end{itemize}

\textbf{Validation target.}
Determine whether leakage remains localized near $|f\rangle$ or spreads into a broader
manifold as predicted by simulation.

\subsection{Experiment F: Filter-Chain Sensitivity}

\textbf{Protocol.}
\begin{itemize}[leftmargin=2em]
    \item Repeat selected power sweeps using different programmed pulse smoothness or
    different effective filtering conditions.
    \item Optionally compare intentionally filtered versus sharp-edge pulses.
\end{itemize}

\textbf{Validation target.}
Test whether the effective transfer-function model in \texttt{cqed\_sim} correctly predicts
shifts in the leakage threshold and changes in non-QND behavior.

\section{Quantitative Validation Metrics}

The report should compare simulation and experiment using explicit metrics:
\begin{enumerate}[leftmargin=2em]
    \item \textbf{Threshold agreement:} difference between simulated and measured power
    where leakage exceeds a chosen benchmark.
    \item \textbf{Curve agreement:} RMSE or weighted fit error for leakage-versus-power and
    leakage-versus-duration curves.
    \item \textbf{Repeated-measurement agreement:} comparison of simulated and measured
    conditional transition probabilities.
    \item \textbf{Distribution agreement:} comparison of means, variances, skewness, and
    tail behavior of integrated I/Q distributions.
    \item \textbf{Regime classification agreement:} whether the simulator correctly
    classifies operating points as approximately QND, moderately leaky, or strongly non-QND.
\end{enumerate}

\section{Software Requirements for \texttt{cqed\_sim}}

The following infrastructure should be built or organized inside \texttt{cqed\_sim}:
\begin{itemize}[leftmargin=2em]
    \item multi-level transmon model interface,
    \item cavity plus transmon Hamiltonian builder,
    \item drive-envelope and transfer-function module,
    \item Lindblad and trajectory solver wrappers,
    \item observable extraction layer,
    \item parameter-sweep and convergence-check framework.
\end{itemize}

Each sweep should automatically generate:
\begin{itemize}[leftmargin=2em]
    \item leakage heatmaps,
    \item ionization heatmaps,
    \item cavity photon traces,
    \item level-population trajectories,
    \item fidelity-versus-leakage tradeoff plots,
    \item synthetic readout histograms,
    \item report-ready tables and figures.
\end{itemize}

\section{Recommended Numerical Campaign}

\subsection{Stage 1: Coarse Survey}
Use a relatively inexpensive model to scan amplitude, pulse length, detuning, $\kappa$, and
$\chi$.

Goal: locate the approximate safe/unsafe boundary.

\subsection{Stage 2: Quantitative Refinement}
Near the boundary, increase transmon-level truncation and cavity truncation, and use more
accurate matrix elements.

Goal: obtain quantitatively stable leakage predictions.

\subsection{Stage 3: Conditioned Record Study}
Select representative points safely below threshold, near threshold, and above threshold.

Goal: generate synthetic readout records and identify experimentally visible signatures of
non-QND behavior.

\subsection{Stage 4: Mitigation Optimization}
Evaluate candidate pulse shapes, detunings, and filter settings at fixed target fidelity.

Goal: identify experimentally implementable protocols that reduce leakage while preserving
speed.

\section{Risks and Caveats}

The report should explicitly note the following caveats:
\begin{itemize}[leftmargin=2em]
    \item High-lying-state ``ionization'' is operational in a truncated model and should
    not be overinterpreted as a literal continuum process.
    \item Inadequate Hilbert-space truncation can produce misleading leakage predictions.
    \item Weakly anharmonic Duffing models may become unreliable at large excitation.
    \item Unknown or poorly calibrated line transfer functions may dominate model-experiment
    discrepancies.
    \item Agreement in assignment fidelity alone is not sufficient to validate the model.
\end{itemize}

\section{Success Criteria}

This project will be considered successful if it produces:
\begin{enumerate}[leftmargin=2em]
    \item a reliable \texttt{cqed\_sim} workflow for simulating measurement-induced leakage
    under strong readout,
    \item regime maps that separate approximately QND from strongly non-QND operation,
    \item level-resolved predictions for $|f\rangle$ and higher-level occupation,
    \item experimentally testable signatures in conditioned readout records,
    \item quantitative comparison procedures between simulation and fridge data,
    \item practical mitigation recommendations for safer fast readout.
\end{enumerate}

\section{Report Structure}

The final report should be organized as follows:
\begin{enumerate}[leftmargin=2em]
    \item Introduction and motivation,
    \item Physical model,
    \item Numerical methods,
    \item Convergence and modeling caveats,
    \item Leakage and ionization regime maps,
    \item Conditioned readout predictions,
    \item Experimental validation methods,
    \item Quantitative simulation-experiment comparison metrics,
    \item Mitigation strategies,
    \item Discussion and next steps.
\end{enumerate}

\section{Agent-Facing Implementation Prompt}

\noindent\textbf{Implementation objective.}
Extend \texttt{cqed\_sim} to support predictive modeling of measurement-induced leakage and
transmon ionization during strong dispersive readout. The framework should support
multi-level transmon plus cavity simulation, unconditional Lindblad evolution, conditioned
quantum trajectories, filtered drive envelopes, and systematic parameter sweeps over
readout amplitude, detuning, $\kappa$, $\chi$, anharmonicity, and pulse shape. The system
should generate level-resolved populations, leakage and high-lying-state metrics, synthetic
readout histograms, repeated-measurement statistics, and report-ready visualizations. The
final report must include explicit experimental validation protocols and quantitative
comparison metrics for fridge data.

\end{document}
